% Part: history
% Chapter: biographies
% Section: rozsa-peter

\documentclass[../../../include/open-logic-section]{subfiles}

\begin{document}

\olfileid{his}{bio}{pet} 

\olsection{روژا پېتر}

\olphoto{peter-rozsa}{روژا پېتر}

روژا پېتر د روژا پوليتسر په نوم د 1905 د
فبرورۍ پر 17مه د هنګري په بوداپېست کښې
زېږېدلې وه۔ هغه تر ټولو زيات په بازګشتي
تابعو د خپل کار له امله پېژندل کېږي،
چې د بازګښت د تيورۍ د ډګر د جوړېدو
دپاره بنسټيز و۔

پېتر په سختو سياسي وختونو کښې لويه شوه؛
په تنکۍ ځوانۍ کښې ئې لومړۍ نړيواله جګړه
روانه وه۔ خو په بوداپېست کښې د ماريا
تېرېزيا د نجونو شتمن ښوونځي ته تللے
شوه او په 1922 کښې ترې فارغه شوه۔ بيا
ئې په بوداپېست کښې د پازمان پېتر په
پوهنتون کښې زده کړه وکړه، چې وروسته
ئې نوم لورانډ اېوتووش پوهنتون شو۔ د
پلار په ټينګار ئې د کيميا زده کړه
پيل کړه، خو وروسته رياضياتو ته
واوښته او په 1927 کښې فارغه شوه۔
که څه هم د ثانوي ښوونځي د رياضياتو
د تدريس اهليت ئې درلود، د هغه وخت
اقتصادي حالت ډېر خراب و، ځکه ستر
اقتصادي رکود د نړۍ اقتصاد اغېزمن
کړے و۔ په دې وخت کښې پېتر د
رياضياتو د مرستندويې ښوونکې او
شخصي ښوونکې په توګه بېلابېل واړۀ
کارونه کول۔ بالاخره پوهنتون ته
ستنه شوه، څو په رياضياتو کښې لوړې
زده کړې وکړي۔ اصل کښې ئې د عددونو
په تيورۍ د کار اراده وه، خو کله
چې پوهه شوه چې نتيجې ئې له مخکې
ثابتې شوې دي، نږدې وه چې رياضيات
بيخي پرېږدي۔ هغه وهڅول شوه چې د
ګوډل د ناتکميلۍ پر قضيو کار وکړي،
او بې له دې چې پوهه وي، د هغه
څو نتيجې ئې په بېلو لارو ثابتې
کړې۔ دې ئې باور بېرته راژوندے
کړ، او پېتر د ډېوېډ هېلبرټ له
بنسټيز پروګرامه په الهام د بازګښت
په تيورۍ خپلې لومړۍ ليکنې وليکلې۔
په 1935 کښې ئې دوکتورا واخيسته،
او په 1937 کښې د \emph{د سمبوليک
منطق مجلې} مديره شوه۔

د پېتر لومړنۍ ليکنې په پراخه توګه د
بازګشتي تابعو د تيورۍ بنسټ ايښودونکې
مرستې ګڼل کېږي۔ په \cite{Peter1935a}
کښې ئې د بازګښت د بېلابېلو ډولونو
اړيکه وڅېړله۔ په \cite{Peter1935b}
کښې ئې وښودله چې يوه ځانګړې په
بازګشتي ډول تعريف شوې تابع بنسټيزه
بازګشتي تابع نۀ ده۔ دې د وېلهېلم
اکېرمان يوه پخوانۍ نتيجه آسانه
کړه۔ د پېتر ساده شوې تابع هغه
ده چې اوس ډېر ځله د اکېرمان
تابع نومېږي، او کله، په لا
کره نوم، د اکېرمان--پېتر تابع۔
هغې د بازګشتي تابعو د تيورۍ
لومړے کتاب وليکۀ \citep{Peter1951}۔

د خپل کار له اهميت او اغېز سره سره پېتر
تر 1945 پورې د بشپړ وخت د تدريس دنده
ترلاسه نۀ کړه۔ د دويمې نړيوالې جګړې
پر وخت د هنګري د نازي اشغال لاندې،
پېتر ته د يهودانو ضد قوانينو له
امله د تدريس اجازه نۀ وه۔ په 1944
کښې حکومت په بوداپېست کښې د
يهودانو يوه محصوره سيمه جوړه کړه؛
دا له پاتې ښاره بېله وه او وسلوالو
ساتونکو ساتله۔ پېتر تر 1945 پورې
هلته ژوند کولو ته اړ وه، کله چې
دا سيمه آزاده شوه۔ بيا ئې د
بوداپېست د ښوونکو د روزنې په
کالج کښې تدريس وکړ، او له 1955
څخه وروسته د اېوتووش لورانډ
په پوهنتون کښې۔ هغه لومړۍ
هنګرۍ رياضيپوهه وه چې د
رياضياتو اکاډميکي دوکتوره شوه،
او لومړۍ ښځه وه چې د هنګري
د علومو اکاډمۍ ته وټاکل شوه۔

پېتر د رياضياتو د شوقمنې ښوونکې په توګه
پېژندل کېده۔ هغې له خپلو زده کوونکو
سره د رياضيکي مسئلو د طبيعت او
ښکلا څېړلو ته تر تش درس ورکولو
ترجيح ورکوله۔ له همدې امله به
زده کوونکو په مينه ورته «روزا
ترور» ويل۔ پېتر په 1977 کښې
د~71 کلونو په عمر وفات شوه۔

\begin{reading}
د نورو ژوندليکي ليکنو دپاره
\citep{Oconnor2014} او \citep{Andrasfai1986}
وګورئ۔ \citet{Tamassy1994} له پېتر سره
يوه لنډه مرکه کړې ده۔ د رياضياتو
د يوه په زړه پورې لوست دپاره د
پېتر کتاب \emph{له نامتناهيت سره
لوبې} وګورئ \citep{Peter2010}۔
\end{reading}

\end{document}
